\subsection{Finite differencing, convergence, error estimates}

We derived the following finite-differencing approximants for the derivative of a function $f(x)$: 
\begin{eqnarray}
\mathit{Forward~differencing:} & f'(x) \simeq \frac{f(x+h)-f(x)}{h} + \mathcal{O}\,(h) \\
\mathit{Backward~differencing:} & f'(x) \simeq \frac{f(x)-f(x-h)}{h} + \mathcal{O}\,(h) \\
\mathit{Central~differencing:} & f'(x) \simeq \frac{f(x+h)-f(x-h)}{2h} + \mathcal{O}\,(h^2)
\end{eqnarray}

\subsubsection*{Problems:}
\begin{enumerate}
\item Write a Python function to compute derivatives using these three finite differencing methods. Compute the derivative of the function $f(x) = e^x \, \sin(x)$ over the range $x = [0, 2\pi)$. Plot the numerically computed derivative $f_{(h)}'(x)$ for three different values of $h$. 
\item Plot the error $\Delta f_{h}'(x) := |f_{h}'(x) - f'(x)|$ for three different values of $h$, and estimate the order of convergence $n$ of each finite-difference approximation. Plot $n(x)$, where 
\begin{equation}
n(x) = \log_2 \frac{f'_{4h}(x) - f'_{2h}(x)}{f'_{2h}(x) - f'_{h}(x)}
\label{eq:order_convg}
\end{equation}
\item Reduce the step size $h$ successively. At what value of $h$ does the round-off error dominate the error budget? 
\item You would notice that the truncation error is not constant across $x$. This calls for non-uniform sampling of $f(x)$ in finite differencing. Derive an explicit central differencing formula assuming that $f(x)$ can be sampled at some \emph{non-uniform} samples $x_i$. 
%\item Derive a central differencing approximant for the first derivative $f'(x)$ that is accurate to $\mathcal{O}(h^4)$. 
\end{enumerate}

\subsection{Richardson extrapolation}

We have seen that, if the order of the error in the numerical estimate of a function $f(x)$ is known, Richardson extrapolation provides a powerful way of improving the accuracy of the estimate. If we have two numerical estimates $f_{h}(x)$ and $f_{2h}(x)$ each having an error of $\mathcal{O}\,(h^k)$, a better estimate is given by 
\begin{equation}
f(x) \simeq \frac{2 ~ 2^k\, f_{h}(x) - f_{2h}(x)}{2\,2^k -1 } + \mathcal{O}\,(h^{l}),
\end{equation}
where $l$ is the next-to-leading-order error term (for e.g., $l = k+2$ for central differencing, while $l = k+1$ for forward/backward differencing). 

\subsubsection*{Problems:}
\label{sec:BBH_nr_data_fd}
\begin{enumerate}
\item Gravitational-waves (GWs) have two independent polarization states -- called ``plus'' and ``cross'' states. GW signals from the coalescence of black-hole binaries, in the simplest case, are circularly polarized: 
\begin{eqnarray}
h_+(t) & = A(t) \, \cos \varphi(t), ~~~ h_\times(t) & = A(t) \, \sin \varphi(t). 
\end{eqnarray}
Download the data file~\cite{nrdata} containing $h_+(t)$ and $h_\times(t)$. (This is the reduced form of the data produced by a numerical-relativity simulation of black-hole binaries performed by the SXS collaboration and is publicly available at~\cite{SXScatalog}). Compute the phase evolution $\varphi(t)$, the frequency evolution $\omega(t) := d\varphi(t)/dt$ and the rate of change of frequency $\dot{\omega}(t) := d\omega(t)/dt$ using second-order central difference approximation. 
\item Estimate the order of convergence of the numerical computation of $\omega(t)$ and $\dot{\omega}(t)$. 
\item Perform an extrapolation of $\omega(t)$ and $\dot{\omega}(t)$ to the next order using estimates of two different time-resolutions.  
\item Derive an explicit expression for $f'(x)$ with error $\mathcal{O}(h^4)$ using Richardson extrapolation. This is the fourth-order finite differencing approximant for the derivative, which we will use later. 

\end{enumerate}
